% 埃米·诺特（综述）
% license CCBYSA3
% type Wiki

本文根据 CC-BY-SA 协议转载翻译自维基百科\href{https://en.wikipedia.org/wiki/Emmy_Noether}{相关文章}

\begin{figure}[ht]
\centering
\includegraphics[width=6cm]{./figures/c0d48f5d45fa6e41.png}
\caption{} \label{fig_AMnt_1}
\end{figure}
阿玛莉·艾米·诺特（Amalie Emmy Noether，[a][b]，1882年3月23日－1935年4月14日）是一位德国数学家，她在抽象代数领域作出了许多重要贡献。她还证明了诺特第一与第二定理，这两项成果在数学物理中具有基础性意义。[4]帕维尔·亚历山德罗夫、阿尔伯特·爱因斯坦、让·迪厄多内、赫尔曼·外尔以及诺伯特·维纳都曾称她为“数学史上最重要的女性”。[5][6][7]作为当时最杰出的数学家之一，诺特发展了环论、域论和代数理论。在物理学中，诺特定理揭示了对称性与守恒定律之间的深刻联系。[8]

诺特出生于法兰克尼亚小镇埃尔朗根的一个犹太家庭，她的父亲是数学家马克斯·诺特。起初，她计划在通过所需考试后教授法语和英语，但后来转而在其父任教的埃尔朗根－纽伦堡大学学习数学。1907年，在保罗·戈尔丹的指导下，她完成了博士论文。此后，她在埃尔朗根数学研究所无薪工作了七年。[9]当时，女性基本上被排除在学术职位之外。1915年，大卫·希尔伯特和费利克斯·克莱因邀请她加入哥廷根大学数学系——这是当时世界知名的数学研究中心。然而哲学学院对此提出反对，因此她在接下来的四年里只能以希尔伯特的名义讲课。1919年，她的任教资格最终获批，从而获得了“私人讲师”的职称。[9]
